% ============================================================
% SST Canon v0.8.1 — New Section 2.16
% Maximal Swirl Tension: Definition, Identities, and
% Threefold Physical Interpretation
% Ready for insertion after Section 2.15 in the main canon
% ============================================================

\subsection{Maximal Swirl Tension}
\label{sec:Fmax}

\subsubsection{Definition and canonical form}

Within the primitive set $\mathcal{P} = \{\rho_f, v_\circlearrowleft, \omega_c\}$,
a characteristic force scale emerges naturally from the Compton energy stored
over the core cross-section. We define the \emph{maximal swirl tension} as

\begin{equation}
  F_{\rm swirl}^{\max}
  \;=\; \frac{v_\circlearrowleft\,\hbar}{2r_c^2}
  \;=\; \frac{\hbar\omega_c}{2r_c}.
  \label{eq:Fmax-def}
\end{equation}

\noindent
[\textsc{Derived}] \textit{Canonical form}: substituting
$v_\circlearrowleft = \omega_c r_c$ into Eq.~\eqref{eq:Fmax-def} gives
$F_{\rm swirl}^{\max} = \hbar\omega_c / 2r_c$,
the Compton energy scale divided by twice the core radius.
Numerically, $F_{\rm swirl}^{\max} = 29.053507\;\text{N}$.

\noindent
[\textsc{Critical Note}] $F_{\rm swirl}^{\max}$ is a mechanical
reformulation of the Compton energy scale over the core cross-section.
It is \emph{not} independent of $\hbar$ and $m_e$; its role in the
canon is interpretive and organisational rather than that of an
independent primitive.

\subsubsection{Internal consistency identities}

The following expressions are algebraically equivalent within the
canonical constant chain and serve as coherence verifications rather
than independent derivations:

\begin{align}
  F_{\rm swirl}^{\max}
  &= \frac{e^2}{16\pi\varepsilon_0 r_c^2},
  \label{eq:Fmax-em}\\
  &= \frac{h\alpha c}{8\pi r_c^2},
  \label{eq:Fmax-hac}\\
  &= \pi r_c^2\,\rho_{\rm calc}\,v_\circlearrowleft^2,
  \label{eq:Fmax-rho}
\end{align}

\noindent
where $\rho_{\rm calc} = 4F_{\rm swirl}^{\max}/(\pi\alpha^2 c^2 r_c^2)$
is the derived bulk density consistent with the canonical core closure.

\noindent
[\textsc{Derived}] \textit{Consistency check}:
Eqs.~\eqref{eq:Fmax-em}--\eqref{eq:Fmax-rho} are equivalent given
$v_\circlearrowleft = \alpha c/2$ and the standard identity
$e^2/(4\pi\varepsilon_0) = \alpha\hbar c$.
Their mutual agreement constitutes a non-trivial numerical
self-consistency test of the primitive set.

\subsubsection{Strongest non-circular identity: the Rydberg constant}

\noindent
[\textsc{Derived}] The Rydberg constant is recoverable directly from
the three SST primitives without $e$, $\varepsilon_0$, or $m_e$ as
explicit input:

\begin{equation}
  R_\infty \;=\; \frac{v_\circlearrowleft^3}{\pi\,r_c\,c^3}.
  \label{eq:Rydberg-SST}
\end{equation}

\noindent
This is the canonical benchmark for the internal non-redundancy of the
primitive set. Combining Eq.~\eqref{eq:Rydberg-SST} with
Eq.~\eqref{eq:Fmax-def} yields the compact relation

\begin{equation}
  \boxed{
  F_{\rm swirl}^{\max}
  \;=\; \frac{32\pi^2\hbar R_\infty^2 c}{\alpha^5},
  }
  \label{eq:Fmax-Rydberg}
\end{equation}

\noindent
which expresses $F_{\rm swirl}^{\max}$ in terms of spectroscopic
observables alone, with no reference to $m_e$, $r_c$, or $e$ as
separate inputs.

\noindent
[\textsc{Derived}] A further identity connecting atomic and
force scales is

\begin{equation}
  R_\infty \;=\; F_{\rm swirl}^{\max}
  \cdot \frac{\alpha^3}{4\pi m_e c^2},
  \label{eq:Rydberg-Fmax}
\end{equation}

\noindent
which states that the Rydberg constant is the maximal swirl tension
scaled by the Compton energy and a purely geometric factor $\alpha^3/4\pi$.
Spectroscopy is thereby connected directly to the æther tension scale.

\subsubsection{Threefold physical interpretation}

The quantity $F_{\rm swirl}^{\max}$ admits three distinct physical
interpretations, each arising from a different sector of physics.
Their numerical coincidence within the canonical constant chain
is a non-trivial coherence result.

\paragraph{(i) Gravitational analogue: the Planck force ceiling.}

In general relativity, the combination $c^4/4G$ sets the maximum
force that any physical process can transmit without forming a
causal horizon~\cite{Gibbons2002}. We denote this scale the maximal
gravitational tension:

\begin{equation}
  F_{\rm gr}^{\max} \;=\; \frac{c^4}{4G}.
  \label{eq:Fgr}
\end{equation}

\noindent
[\textsc{Orthodox}] $F_{\rm gr}^{\max}$ is the Planck force,
an invariant upper bound in classical general relativity.

\noindent
[\textsc{Derived}] The ratio of the two tension scales reduces to

\begin{equation}
  \boxed{
  \frac{F_{\rm swirl}^{\max}}{F_{\rm gr}^{\max}}
  \;=\; \frac{4\alpha_g}{\alpha},
  }
  \label{eq:ratio}
\end{equation}

\noindent
where $\alpha_g = Gm_e^2/\hbar c$ is the gravitational
fine-structure constant of the electron and $\alpha$ is the
electromagnetic fine-structure constant.
Numerically, $4\alpha_g/\alpha \approx 9.6\times10^{-43}$,
in agreement with the directly computed ratio
$29.053507\,\text{N} / 3.026\times10^{43}\,\text{N}$.

\noindent
[\textsc{Speculative}] \textit{Interpretation}:
the hierarchy between the electromagnetic and gravitational
force scales in SST is not an unexplained numerical accident
but the ratio of the two fundamental coupling constants of the
electron. The medium is $\sim\!10^{42}$ times stiffer than the
space-time geometry, and this stiffness ratio is exactly
$4\alpha_g/\alpha$.

\paragraph{(ii) Electrostatic interpretation: the Coulomb ceiling
at core scale.}

When two protons are compressed to a centre-to-centre separation
equal to the vortex core radius $r_c$, the Coulomb repulsion is

\begin{equation}
  F_{pp}(r_c)
  \;=\; \frac{e^2}{4\pi\varepsilon_0 r_c^2}
  \;=\; 4\,F_{\rm swirl}^{\max}.
  \label{eq:Coulomb-core}
\end{equation}

\noindent
[\textsc{Derived}] The maximal Coulomb repulsion on the core scale
is exactly four times $F_{\rm swirl}^{\max}$.
The factor of four matches the denominator structure of
Eq.~\eqref{eq:Fgr}, and is not numerically coincidental within
the canonical chain.

\noindent
[\textsc{Speculative}] \textit{Interpretation}: $F_{\rm swirl}^{\max}$
is the quarter-Coulomb force at core scale. The medium can sustain
one quarter of the maximum proton--proton Coulomb repulsion before
the topological structure of the vortex core fails. This provides a
physical ceiling on the compression of charged vortex configurations.

\paragraph{(iii) Mechanical interpretation: the spring constant of
photon--electron interaction.}

The VAM appendix~\cite{VAMAppendix} models the electron's internal
vortex as a linear spring with stiffness

\begin{equation}
  K_e \;=\; \frac{F_{\rm swirl}^{\max}}{n\,r_c},
  \qquad n = 2,
  \label{eq:spring-K}
\end{equation}

\noindent
where $n$ counts the number of half-wavelength segments on the core
and is fixed empirically to $n = 2$ by the photon--electron matching
condition. The natural oscillation frequency of the electron core
under this spring is

\begin{equation}
  \omega_{\rm spring}
  \;=\; \sqrt{\frac{K_e}{m_e}}
  \;=\; \sqrt{\frac{F_{\rm swirl}^{\max}}{2\,r_c\,m_e}}.
  \label{eq:omega-spring}
\end{equation}

\noindent
Substituting Eq.~\eqref{eq:Fmax-def} and $m_e = \hbar\omega_c/c^2$:

\begin{equation}
  \omega_{\rm spring}
  \;=\; \frac{\omega_c}{\alpha}
  \;=\; \frac{m_e c^2}{\alpha\hbar},
  \label{eq:omega-spring-result}
\end{equation}

\noindent
a frequency intermediate between the atomic Rydberg scale
$(\sim\!\alpha^2\omega_c)$ and the Compton scale $(\omega_c)$.

\noindent
[\textsc{Derived}] The energy stored in this spring at maximum
compression $\delta = r_c$ is

\begin{equation}
  E_{\rm spring}
  \;=\; \tfrac{1}{2}\,K_e\,r_c^2
  \;=\; \frac{F_{\rm swirl}^{\max}\,r_c}{4}
  \;=\; \frac{\hbar\omega_c}{8}
  \;=\; \frac{m_e c^2}{8}.
  \label{eq:Espring}
\end{equation}

\noindent
One eighth of the electron rest energy is stored as internal
vortex-spring potential at maximum compression.

\noindent
[\textsc{Speculative}] \textit{Physical picture}: during
photon absorption, the electron's vortex core is compressed
for a duration of order one Planck time $t_p$ into a transient
mass-like configuration. The restoring force scale for this
compression is $F_{\rm swirl}^{\max}$, and the energy scale
of the transient is $m_e c^2/8$, one eighth of the rest energy.
The spring releases at the Compton period, generating the
recoil that initiates the Coulomb interaction with the proton.

\noindent
[\textsc{Critical Note}] Equation~\eqref{eq:omega-spring-result}
assumes $n = 2$ from empirical matching. A first-principles
derivation of the photon wrap number from the delay dynamics
of Section~\ref{sec:delay} is an open step.

\subsubsection{Unified summary}

The three interpretations are collected in Table~\ref{tab:Fmax}.
Together they identify $F_{\rm swirl}^{\max}$ as the
\emph{structural tension threshold of the medium}: the scale at
which a physical system --- gravitational, electrostatic, or
mechanical --- exhausts its capacity to store energy in the
current configuration and must either fail topologically,
form a horizon, or release energy as radiation.

\begin{table}[h]
\centering
\caption{Threefold interpretation of $F_{\rm swirl}^{\max}$
         within the SST canonical framework.}
\label{tab:Fmax}
\begin{tabular}{llll}
\hline
\textbf{Context} & \textbf{Expression} & \textbf{Physical meaning}
  & \textbf{Status} \\
\hline
Gravitational ceiling &
  $F_{\rm gr}^{\max}/F_{\rm swirl}^{\max} = \alpha/(4\alpha_g)$ &
  Medium stiffer than space-time by $\alpha/4\alpha_g$ &
  [Derived] \\
Coulomb ceiling &
  $F_{pp}(r_c) = 4\,F_{\rm swirl}^{\max}$ &
  Quarter of max.\ Coulomb force at core scale &
  [Derived] \\
Vortex spring &
  $K_e = F_{\rm swirl}^{\max}/2r_c$ &
  Spring constant of photon--electron interaction &
  [Speculative] \\
\hline
\end{tabular}
\end{table}

\noindent
[\textsc{Speculative}] \textit{Canonical interpretation}:
$F_{\rm swirl}^{\max}$ is the tension at which the structured
æther medium reaches its topological integrity limit. It plays
the same organisational role for the electromagnetic and
mechanical sectors that $F_{\rm gr}^{\max} = c^4/4G$ plays for
the gravitational sector. Their ratio is the ratio of the two
fundamental coupling constants of the electron.

% ============================================================
% References to add to main bibliography:
%
% \bibitem{Gibbons2002}
% G.~W.~Gibbons,
% The maximum tension principle in general relativity,
% Found.\ Phys.\ \textbf{32} (2002) 1891--1901,
% DOI: 10.1023/A:1022370717626.
%
% \bibitem{VAMAppendix}
% O.~Iskandarani,
% Appendix X: From \AE{}ther Tension to Planck's Constant
% and the Bohr Radius (internal SST working document, 2025).
% ============================================================
